\section{Séparer déclaration et définition : les fichiers d'en-tête}

\subsection{Le problème}

Quand un programme grossit, il devient impossible (et peu pratique) de tout écrire dans un seul fichier \texttt{.c}. On souhaite pouvoir :
\begin{itemize}
  \item regrouper des fonctions qui vont ensemble dans leur propre fichier (par exemple, toutes les fonctions de conversion de température) ;
  \item réutiliser ces fonctions dans plusieurs programmes différents, sans recopier leur code ;
  \item permettre à un autre fichier \texttt{.c} d'utiliser ces fonctions sans avoir à en connaître le détail d'implémentation.
\end{itemize}

Pour cela, le C sépare deux choses :
\begin{description}
  \item[la déclaration] (le prototype) : \emph{ce que fait} la fonction (son nom, ses paramètres, son type de retour), placée dans un fichier d'en-tête \texttt{.h} ;
  \item[la définition] (le corps) : \emph{comment} elle le fait, placée dans un fichier source \texttt{.c}.
\end{description}

\begin{UPSTIinfor}{Rôle d'un fichier d'en-tête (A retenir !)}
  Un fichier \texttt{.h} contient uniquement les \textbf{prototypes} des fonctions (et éventuellement des constantes, des types). Il ne contient \textbf{pas} le corps des fonctions. Il sert de « carte de visite » : il décrit ce qu'un fichier \texttt{.c} propose, sans en dévoiler le contenu.
\end{UPSTIinfor}

\subsection{Exemple : un mini-module de géométrie}

On souhaite regrouper des fonctions de calcul géométrique dans un module réutilisable, composé de deux fichiers.

\begin{lstlisting}[language=c,caption={\texttt{geometrie.h} : les déclarations}]
#ifndef GEOMETRIE_H
#define GEOMETRIE_H

float aire_disque(float rayon);
float aire_triangle(float base, float hauteur);

#endif
\end{lstlisting}

\begin{lstlisting}[language=c,caption={\texttt{geometrie.c} : les définitions}]
#include "geometrie.h"

float aire_disque(float rayon) {
    return 3.14159 * rayon * rayon;
}

float aire_triangle(float base, float hauteur) {
    return base * hauteur / 2.0;
}
\end{lstlisting}

\begin{lstlisting}[language=c,caption={\texttt{main.c} : utilisation du module}]
#include <stdio.h>
#include "geometrie.h"

int main(void) {
    printf("Aire du disque : %.2f\n", aire_disque(2.0));
    printf("Aire du triangle : %.2f\n", aire_triangle(4.0, 3.0));
    return 0;
}
\end{lstlisting}

Le fichier \texttt{main.c} peut utiliser \texttt{aire\_disque} et \texttt{aire\_triangle} sans jamais avoir vu leur corps : il lui suffit de connaître leur prototype, fourni par \texttt{geometrie.h}.

\begin{UPSTIwarning}{\texttt{\#include <...>} ou  \texttt{\#include "..."} ?}
  \begin{itemize}
    \item \texttt{\#include <stdio.h>} : chevrons, pour une bibliothèque \textbf{standard} du langage (fournie par le compilateur).
    \item \texttt{\#include "geometrie.h"} : guillemets, pour un fichier d'en-tête \textbf{du projet}, situé (en général) dans le même dossier que le fichier source.
  \end{itemize}
\end{UPSTIwarning}

\subsection{La protection contre les inclusions multiples}

Si un fichier d'en-tête est inclus deux fois dans le même programme (ce qui arrive facilement quand plusieurs fichiers \texttt{.c} incluent le même \texttt{.h}), le compilateur verrait deux fois les mêmes prototypes, ce qui provoque des erreurs de compilation.

\begin{UPSTIinfor}{Le garde d'inclusion (\emph{include guard})}
  On protège chaque fichier \texttt{.h} avec les trois lignes suivantes :
  \begin{lstlisting}[language=c,caption={Garde d'inclusion}]
#ifndef GEOMETRIE_H   // si GEOMETRIE_H n'est pas deja defini
#define GEOMETRIE_H   // on le definit (pour la prochaine fois)

// ... contenu du fichier .h ...

#endif                // fin de la protection
  \end{lstlisting}
  À la première inclusion, \texttt{GEOMETRIE\_H} est défini et le contenu est lu normalement. À une éventuelle deuxième inclusion, \texttt{GEOMETRIE\_H} est déjà défini : le contenu entre \texttt{\#ifndef} et \texttt{\#endif} est ignoré.

  Une alternative, plus courte mais moins portable, consiste à écrire simplement \texttt{\#pragma once} en première ligne du fichier \texttt{.h}.
\end{UPSTIinfor}

\begin{UPSTIactivite}[][Écrire un fichier d'en-tête]
On souhaite créer un module \texttt{conversions} regroupant deux fonctions : \texttt{float km\_vers\_miles(float km)} et \texttt{float miles\_vers\_km(float miles)}.

\UPSTIquestion{Écrire le contenu du fichier \texttt{conversions.h}, avec sa protection contre les inclusions multiples.}
\UPSTIquadrillage{3}
\UPSTIquestion{Quelle ligne faut-il ajouter en tête de \texttt{conversions.c} pour qu'il ait accès à ses propres prototypes ?}
\UPSTIquadrillage{1}
\end{UPSTIactivite}

\begin{UPSTIprofOnlyEnv}
\begin{UPSTIcorrectionP}{Écrire un fichier d'en-tête}
\begin{lstlisting}[language=c,caption={\texttt{conversions.h}}]
#ifndef CONVERSIONS_H
#define CONVERSIONS_H

float km_vers_miles(float km);
float miles_vers_km(float miles);

#endif
\end{lstlisting}
Il faut ajouter \texttt{\#include "conversions.h"} en tête de \texttt{conversions.c} : cela permet au compilateur de vérifier que les fonctions définies correspondent bien aux prototypes annoncés.
\end{UPSTIcorrectionP}
\end{UPSTIprofOnlyEnv}

\subsection{Compiler plusieurs fichiers \texttt{.c} ensemble}

Un programme peut être réparti sur plusieurs fichiers \texttt{.c}, chacun avec son fichier \texttt{.h} associé. Il faut alors indiquer au compilateur \textbf{tous} les fichiers \texttt{.c} à compiler ensemble :

\begin{lstlisting}[language=bash,caption={Compilation de plusieurs fichiers sources}]
gcc main.c geometrie.c -o programme
\end{lstlisting}

\begin{UPSTIinfor}{Ce qu'il faut retenir sur les fichiers d'en-tête}
  \begin{itemize}
    \item Un \texttt{.h} contient des \textbf{prototypes} (déclarations), un \texttt{.c} contient le \textbf{corps} des fonctions (définitions).
    \item On protège chaque \texttt{.h} avec un garde d'inclusion (\texttt{\#ifndef}/\texttt{\#define}/\texttt{\#endif} ou \texttt{\#pragma once}).
    \item Tous les fichiers \texttt{.c} utilisés par un programme doivent être passés au compilateur en même temps.
  \end{itemize}
\end{UPSTIinfor}

Pour des projets plus gros, comportant de nombreux fichiers, on utilise en général un outil qui automatise cette compilation (un \texttt{Makefile}, par exemple) : ce sera vu plus tard.
